\section{Evaluating current knowledge on the dynamic stability of tastes}
In which direction does the extant evidence point?  While there has been little empirical research on the dynamic stability of tastes through time, there is good reason to suppose that they are more stable than the network-transmission model indirectly implies.  Consider the following well-established findings:

\begin{enumerate}
\item  Most studies of the role of early family and school experiences on cultural capital have shown strong effects of arts participation, education and after-school training during adolescence on adult tastes even after controlling for subsequent educational attainment \citep{kracman96}.\marginpar{Talk about DiMaggio's 1982 paper} Hagan's \citeyearpar{hagan91} classic study of the effect of subcultural preferences on adult stratification outcomes implicitly supports a stability assumption. Hagan finds that involvement in a ``party subculture'' in adolescence has deleterious effects on the career advances of working class youth, but actually \emph{helps} the attainment chances of middle class youth, once the negative effects of these set of subcultural preferences on school performance is adjusted for. Hagan's concludes that it is precisely because the cultural habits obtained during this episode of ``drift'' continue to operate in adulthood that we find this class-specific enabling effect.

\item Studies that look at the patterning of taste across age groups show that specific genre preferences (e.g. for musical styles or even particular songs) are acquired early in adolescence and do not change dramatically afterwards. \citet{smith95} using data from the 1993 General Social Survey, shows that musical tastes are developed early in young adulthood and appear to be fairly distinct across age-cohorts.\marginpar{add quote from Smith and Holbrook}  

\item The few studies that are able to speak to the stability issue, indeed find support for the cultural-capital implication.  \citep[178]{bennett_etal99}, using qualitative evidence from retrospective interview reports on musical preferences in Australia find that ``{\ldots}many of those who were in middle age had retained  the music preference of their late teens and early adulthood.'' They conclude that while they ``{\ldots}cannot discount the possibility that some people will undergo radical changes in their tastes as they age, on bzalance the evidence appears to favor the explanation that the taste acquired in the formative years of a cohort persist into later life'' \citep[180]{bennett_etal99}.  In a recent study in Germany \citet{georg04}  reinterviewed a sample of respondents in 2002 which had previously been interviewed in 1983.  In the initial study, which was conducted when most respondents were adolescents, respondents provided information on a series of items regarding the participants preferences for various form of high-status cultural activities (e.g. attending a classical music or jazz concert; reading serious literature, etc.), most of which were then replicated in 2002. Drawing on the cultural-capital expectation of taste stability, Georg hypothesized that cultural preferences for high-status culture should be relatively stable throughout the life course. Using cross-lagged structural equations models he finds striking support this hypothesis ($\beta=0.61$, with $\beta=1.0$ being perfect stability and $\beta=0$ being no-stability).
\end{enumerate}•

Taken together, these findings suggest that while the effect of social network change on cultural taste loss is not denied, whatever relational (or ``network'') effects on taste exist, must operate in the background of high-levels of cross-temporal taste stability and the reciprocal action of cultural taste on social networks. Thus, contrary to the network transmission account, and in accordance with the imagery proposed in the culture-conversion model, tastes (in contrast to network ties) do not exhibit over-time volatility, but instead show high inter-temporal correlation and ``stickiness'' in the following section I explore some important implications of this conclusion. 
